\section{Brechas identificadas}
A pesar de los avances en biotecnología y digitalización, persisten importantes brechas que limitan la implementación efectiva de soluciones sostenibles para la gestión de residuos agrícolas y la mitigación de gases de efecto invernadero (GEI). Estas brechas incluyen la falta de sistemas integrados para el monitoreo y control de variables ambientales, la escasa aplicación de tecnologías como IoT e inteligencia artificial, y la optimización deficiente en procesos de bioconversión con mosca soldado negra. Además, se identifican limitaciones en la accesibilidad y escalabilidad de soluciones tecnológicas, así como una débil integración entre biotecnología y herramientas digitales, lo que refuerza la necesidad de desarrollar enfoques innovadores que superen estas barreras.

\begin{itemize}
    \item Falta de sistemas integrados para el monitoreo y control de variables ambientales.
    \item Limitaciones en la medición de gases de efecto invernadero (GEI).
    \item Escasa aplicación de IoT e IA en la gestión de residuos agrícolas.
    \item Optimización deficiente en la biodegradación con Mosca Soldado Negra.
    \item Baja accesibilidad y escalabilidad de soluciones tecnológicas.
    \item Débil integración entre biotecnología y tecnologías digitales.
    \item Insuficiencia de datos estandarizados y modelos predictivos.
    \item Impacto limitado en la mitigación del cambio climático.
\end{itemize}

Estas brechas justifican la necesidad de implementar soluciones innovadoras que integren IoT e IA para mejorar la gestión de residuos agrícolas, optimizar los procesos de medición, reporte y validación de GEI, y avanzar hacia un modelo de agricultura sostenible y climáticamente responsable.